% !TEX root = ../Thesis.tex

\chapter{Stato dell'arte}\label{ch:stato-arte}

\section[Meteor]{Meteor: un framework di steganografia generativa su testo}

\subsection{Contesto e motivazioni}
La necessità di una solida comprensione teorica, indispensabile per contestualizzare appieno il presente lavoro di ricerca, si articola attraverso una costellazione di concetti interdipendenti che convergono in un'unica opera fondamentale: il lavoro di \cite{kaptchuk-2021}, comunemente noto come \emph{Meteor}. Questa pionieristica indagine, che va oltre una mera disamina della steganografia classica e di quella generativa, introduce primariamente un approccio metodologico e, in secondo luogo, un framework robusto ed efficiente deputato alla \emph{generazione di messaggi testuali contenenti informazioni celate}.

La peculiarità di queste informazioni risiede nella loro intrinseca invisibilità non solo all'occhio umano, ma anche, e soprattutto, alle più sofisticate tecniche di steganalisi basate su divergenza statistica. La ragione di tale elusività risiede nella formalizzazione del concetto di \emph{perfectly secure steganography} nel dominio del linguaggio naturale: il testo generato da \emph{Meteor} non manifesta alcuna deviazione statistica significativa rispetto a un testo prodotto autonomamente dal medesimo modello linguistico di base (il cosiddetto \emph{cover distribution}). Come verrà dettagliato in seguito, \emph{Meteor} si avvale di un modello linguistico che ha raggiunto una notevole notorietà, estendendosi ben oltre gli ambiti accademici tradizionali, configurandosi come uno dei primi Transformer a larga scala a influenzare il dibattito pubblico sulla capacità generativa delle IA. Nello specifico, \emph{Meteor} impiega il modello \textbf{GPT-2} (Generative Pre-trained Transformer 2), caratterizzato da un'architettura \emph{decoder-only} addestrata su un corpus massivo di dati testuali (WebText). Tuttavia, l'architettura del framework è stata concepita con una notevole flessibilità, risultando agnostica rispetto al modello linguistico sottostante (\emph{Backbone LLM}). Questa modularità permette a \emph{Meteor} di beneficiare direttamente dei progressi nello stato dell'arte dei modelli linguistici, poiché l'efficacia della steganografia è intrinsecamente legata alla capacità del modello di approssimare la vera distribuzione del linguaggio umano ($P_{data}$).

Per una comprensione ancora più profonda, è opportuno esaminare il contesto socio-tecnologico in cui \emph{Meteor} è maturato e le problematiche specifiche che si proponeva di risolvere, con particolare attenzione alla crescente sofisticazione dei sistemi di sorveglianza e censura digitale. In un'epoca in cui la libertà di espressione e la privacy online sono costantemente minacciate da regimi autoritari e da sistemi di filtraggio sempre più avanzati, la necessità di strumenti di comunicazione \emph{censorship-resistant} è diventata più urgente che mai. \emph{Meteor} si inserisce in questo contesto come una risposta innovativa, offrendo un metodo per nascondere informazioni all'interno di testi generati in modo naturale, rendendo estremamente difficile per qualsiasi osservatore esterno rilevare la presenza di un messaggio segreto.

L'essenza di \emph{Meteor} risiede nell'idea innovativa di utilizzare un modello generativo per la creazione di contenuti steganografici, sfruttando la sua capacità di modellare la probabilità condizionata $P(x_t | x_{<t})$. In termini pratici, il modello agisce come un sofisticato stimatore di densità: data una sequenza di contesto iniziale (\emph{prompt}), esso fornisce una distribuzione di probabilità dettagliata (solitamente tramite una funzione \emph{softmax} applicata ai \emph{logits}) per tutti i possibili \emph{token} contenuti nel vocabolario $\mathcal{V}$. La scelta del \emph{token} da inserire nella sequenza viene effettuata tramite un meccanismo di campionamento che mappa i bit del messaggio segreto in indici del vocabolario, seguendo l'ordinamento indotto dalla distribuzione di probabilità. Questa metodologia trasforma il processo di inferenza in un'operazione di codifica di canale: il testo risultante è statisticamente probabile secondo il modello, ma contiene un messaggio segreto decodificabile univocamente dal destinatario. La decodifica è resa possibile a condizione che il ricevente disponga dello stesso modello linguistico, della medesima sequenza di contesto e di una chiave privata condivisa (\emph{shared secret}), necessaria per inizializzare il generatore di numeri pseudo-casuali o per definire le permutazioni degli indici del vocabolario.

Costruiamo ora uno scenario illustrativo volto a consolidare la comprensione del meccanismo operativo di \emph{Meteor}. Immaginiamo una comunicazione confidenziale tra due entità, Alice ($A$) e Bob ($B$), che desiderano scambiarsi informazioni sensibili attraverso un canale pubblico monitorato. L'obiettivo primario è impedire a un osservatore avversario, Eve ($E$), di rilevare la presenza stessa della comunicazione (\emph{steganographic undetectability}). Un contesto critico si verifica quando Alice o Bob operano all'interno di infrastrutture soggette a \emph{Deep Packet Inspection} (DPI) o filtraggio governativo stringente. In tali scenari, i messaggi cifrati tramite crittografia standard (come AES o RSA) presentano un'elevata entropia e caratteristiche distintive che li rendono facilmente identificabili come traffico cifrato, risultando facilmente identificabili come traffico non lecito e quindi soggetto a blocco immediato (\autoref{fig:alice-bob-eve-firewall}). Un esempio emblematico è il sistema descritto in \autocite{chinawall-2015}, il "\emph{Great Firewall}", che impiega euristiche avanzate per discriminare tra traffico legittimo e protocolli di tunneling o cifratura.
All'interno di questo quadro ostile, la steganografia generativa supera il limite della crittografia tradizionale: non si limita a proteggere il contenuto, ma nasconde l'esistenza stessa dell'atto comunicativo. Eve, osservando il transito di un testo coerente e stilisticamente appropriato, non ha basi statistiche per etichettare la comunicazione come sospetta\footnote{Nel presente scenario, "segreto" e "privato" sono usati in senso tecnico: una comunicazione è steganograficamente sicura se la divergenza di Kullback-Leibler tra la distribuzione dei messaggi naturali e quella dei messaggi steganografici tende a zero.}.

\begin{figure}[tbh]
	\centering
	\frame{\includegraphics[width=.8\textwidth]{AliceBobEveFirewall.pdf}}
	\caption[Alice e Bob attraversano il firewall di Eve]{Rappresentazione schematica di una comunicazione tra Alice (A) e Bob (B) in presenza di un osservatore esterno, Eve (E), e di un sistema di filtraggio dei contenuti (Firewall) che analizza la firma statistica dei pacchetti.}
	\label{fig:alice-bob-eve-firewall}
\end{figure}

Sotto il profilo teorico, la progettazione di \emph{Meteor} ha richiesto la risoluzione di due criticità fondamentali intrinseche alla natura del linguaggio naturale.

La prima sfida concerne la definizione di un \emph{sampler steganografico} efficace. L'algoritmo deve mappare il contesto fornito al modello in una distribuzione di probabilità del token successivo che ricalchi quanto più fedelmente possibile quella del linguaggio naturale. Risulta tuttavia evidente che il linguaggio naturale non sottostia a una distribuzione di probabilità assoluta e invariante: si considerino, a titolo esemplificativo, le divergenze linguistiche intercorrenti tra un testo dei primi del Novecento e uno contemporaneo, o tra la prosa scientifica e quella letteraria. \emph{Meteor} affronta tale complessità adottando un modello di generazione testuale come approssimatore dinamico del linguaggio naturale. Tale approccio consente di simulare efficacemente la distribuzione linguistica in molteplici domini, superando il problema dell'assenza di una distribuzione statistica universale che ha storicamente limitato la steganografia classica. È tuttavia opportuno precisare che l'efficacia di \emph{Meteor} è intrinsecamente dipendente dalla capacità del modello linguistico sottostante (il \emph{Backbone LLM}) di approssimare la distribuzione reale. Qualora il modello non riuscisse a catturare le sfumature e le complessità semantiche del linguaggio umano, la steganografia generativa diverrebbe vulnerabile a tecniche di rilevamento avanzate. Cionondimeno, ogni output prodotto da \emph{Meteor} risulta statisticamente indistinguibile da un testo generato autonomamente dal medesimo modello; ciò garantisce una sicurezza steganografica assoluta, a condizione che non sia possibile distinguere tra un testo prodotto dal modello e uno naturale.

La seconda sfida è legata alla variabilità dell'entropia testuale. In segmenti caratterizzati da elevata densità semantica o vincoli sintattici stringenti (bassa entropia), la libertà di selezione del token successivo è pressoché nulla. Ad esempio, data la sequenza \emph{"The largest carnivore of the Cretaceous period was the Tyrannosaurus"}, il "token" \emph{"Rex"} presenta una probabilità estremamente elevata. Forzare la codifica di bit in tale posizione selezionando un token alternativo produrrebbe un'anomalia semantica macroscopica, facilmente individuabile da un classificatore neurale. Per ovviare a ciò, \emph{Meteor} introduce una soglia di entropia minima: il sistema analizza la distribuzione di probabilità a ogni passo temporale e determina dinamicamente se la varianza sia sufficiente per ospitare informazione. Se l'entropia è inferiore a un valore critico $\tau$, il sistema genera un token senza inserire informazione steganografica. Qualora, invece, l'entropia superi tale soglia, il sistema seleziona i token con probabilità rilevante, escludendo la coda della distribuzione. Attraverso un meccanismo di partizionamento degli indici basato sulle probabilità relative, viene campionato il token i cui estremi dell'intervallo di probabilità presentino una corrispondenza tra i bit più significativi e i bit del messaggio da codificare. Nel caso in cui l'intervallo non permetta la codifica di alcun bit (ovvero quando gli estremi non condividono il prefisso binario), \emph{Meteor} adotta una strategia a bitrate variabile, rinunciando all'inserimento di informazione per quel determinato passo. In fase di decodifica, la conoscenza della distribuzione di probabilità e dei relativi indici consente di determinare univocamente se l'informazione sia stata inserita o meno, potendo recuperare l'intervallo relativo al token selezionato e recuperando i bit più significativi comuni per tutti i valori nell'intervallo. Sebbene tale eventualità sia rara data l'elevata entropia media del linguaggio, questa strategia di codifica adattiva (\emph{adaptive rate}) garantisce che il throughput informativo sia massimizzato esclusivamente dove il linguaggio offre una ridondanza naturale, preservando rigorosamente l'integrità statistica del testo prodotto.

\subsection{Meccanismo Operativo del Protocollo Meteor}

Il concetto fondamentale che differenzia Meteor dai metodi tradizionali di steganografia è lo spostamento del paradigma: l'informazione non viene "nascosta" all'interno di un oggetto esistente tramite modifiche (embedding), ma viene utilizzata per \textbf{guidare il processo di campionamento} durante la generazione di un nuovo media.

Il funzionamento pratico del protocollo si articola nelle seguenti fasi sequenziali:

\begin{enumerate}
	\item \textbf{Cifratura del Messaggio:} Il messaggio segreto viene cifrato utilizzando una chiave simmetrica condivisa tra mittente e destinatario. Questo passaggio è essenziale per trasformare il messaggio in una stringa di bit computazionalmente indistinguibile dal rumore casuale, garantendo che le scelte di campionamento effettuate dal protocollo non presentino pattern statistici rilevabili.

	\item \textbf{Accesso alla Distribuzione Condizionata:} Meteor si appoggia a un modello generativo pre-addestrato (originariamente GPT-2 il quale possiede un'architettura basata su Transformers, necessità imprescindibile anche per lo sviluppo di modelli nel dominio visivo, quindi anche nel caso della VQ-GAN utilizzata nella nostra implementazione, che vedremo in seguito). Ad ogni passo temporale $t$, il modello riceve il contesto dei token già generati e produce una \textbf{distribuzione di probabilità $P$} su tutto il vocabolario dei possibili token successivi.

	\item \textbf{Campionamento Vincolato:} Invece di campionare casualmente o scegliere il token più probabile, Meteor utilizza i bit del messaggio cifrato come "indice" per selezionare il token dalla distribuzione $P$. Lo spazio di probabilità viene partizionato e mappato sulle sequenze di bit in modo che il token scelto sia sempre \textbf{statisticamente coerente} con il modello, selezionando il token tra quelli più probabili, eliminando la coda della distribuzione sotto una soglia pre-specificata di probabilità. Poiché la scelta avviene all'interno della distribuzione naturale del modello, l'output steganografico rimane indistinguibile da un output generato classicamente dallo stesso modello.

	\item \textbf{Generazione Iterativa:} Questo processo avviene in modo autoregressivo. Ogni token generato entra a far parte del nuovo contesto per il passo successivo, continuando fino a quando l'intero media (o la sequenza di token visivi necessari per ricostruire l'immagine tramite il decoder VQ-GAN) non è completo. Il troncamento non è netto dopo la codifica dell'ultimo token, ma può essere definito da un carattere speciale che indica la fine del messaggio, o da un limite di lunghezza predefinito. In questo modo, il testo o l'immagine risultante è completamente naturale e coerente, senza alcuna evidenza della presenza di un messaggio nascosto.

	\item \textbf{Estrazione e Decodifica:} Il ricevente, possedendo lo stesso modello generativo e la chiave simmetrica, è in grado di ricreare le medesime distribuzioni di probabilità $P$ per ogni passo del processo. Identificando quali token sono stati effettivamente inclusi nello stego-media ricevuto, il destinatario può risalire ai bit che hanno determinato quelle scelte, ricostruendo il messaggio cifrato e, infine, il segreto originale.
\end{enumerate}

L'efficacia di Meteor risiede nella sua capacità di eliminare le anomalie statistiche che gli strumenti di \textbf{steganalisi moderna} solitamente identificano. In questo contesto, il messaggio segreto non è contenuto "nei" dati, ma rappresenta la \textbf{ragione deterministica} per cui determinati dati coerenti sono stati generati al posto di altri.

Per comprendere appieno il flusso del protocollo, è utile illustrare uno scenario concreto\footnote{L'esempio proposto rappresenta una forzatura narrativa volta a evitare i classici scenari di sorveglianza bellica o regimi autoritari. Sebbene la trama possa apparire singolare, ogni passaggio tecnico descritto riflette operazioni realmente attuabili a livello pratico.}. Supponiamo che Alice voglia comunicare con il compagno Bob senza che sua madre Eve, esperta di reti di telecomunicazioni e ostile alla loro relazione, lo venga a sapere. Alice ha dimenticato il proprio computer a casa di Bob e deve chiedergli di riportarglielo, ma Eve intercetta ogni comunicazione all'interno della rete domestica. In questo scenario, Alice ha esaurito il traffico dati cellulare e non ha accesso a reti Wi-Fi alternative. Il firewall di casa è stato configurato da Eve per bloccare qualsiasi traffico crittografato end-to-end; i messaggi devono transitare in chiaro per poter superare il gateway. La navigazione esterna è mediata da un proxy che instrada le informazioni tramite un protocollo, ideato dalla stessa Eve, che lascia i contenuti leggibili. Di conseguenza, ogni messaggio inviato o ricevuto da Alice è potenzialmente visibile a sua madre.

In questo contesto, qualsiasi comunicazione cifrata in forma tradizionale verrebbe immediatamente bloccata. Alice e Bob decidono quindi di utilizzare \emph{Meteor}. Alice redige un messaggio segreto: \enquote{Mi porteresti il portatile domani per favore, l'ho lasciato sul tavolo di cucina da te.}. Utilizzando una chiave simmetrica condivisa (ad esempio la data del loro anniversario), Alice cifra il testo ottenendo una stringa di bit. Il messaggio originale, lungo 87 caratteri in codifica ASCII, corrisponde a 87 byte, ovvero 696 bit. Dopo la cifratura, si ottiene una sequenza di 696 bit che è computazionalmente indistingubile da una stringa interamente casuale. Poiché l'invio diretto di tale stringa verrebbe intercettato e bloccato da Eve, Alice si affida a \emph{Meteor} per trasformarla.

Eve è abituata a vedere Alice e Bob scambiarsi testi generati da modelli linguistici mentre studiano insieme per l'università. Alice, studentessa di ingegneria informatica, decide di sfruttare questa abitudine: sa che Eve non sospetterà nulla se vedrà transitare nella chat una serie di messaggi relativi alla preparazione dell'esame di Analisi II, scambiandoli per normali interazioni con un LLM. Alice inserisce il messaggio cifrato in \emph{Meteor} e fornisce come contesto iniziale (\emph{prompt}) il testo di un esercizio sullo studio del dominio di una funzione. Il framework genererà un testo che non solo nasconde il messaggio cifrato in modo trasparente, ma risulterà anche semanticamente coerente con il contesto matematico fornito.

L'efficienza di questo processo dipende dalla \enquote{precisione} di \emph{Meteor}, ovvero dal numero di bit che possono essere codificati in un singolo token. Questo valore è influenzato dalla dimensione del vocabolario e dalla distribuzione di probabilità dei token. Mentre per modelli come GPT-2 è teoricamente possibile codificare fino a 32 bit per token, in media se ne ottengono circa 15. Tuttavia, in contesti a bassa entropia come un esercizio di analisi, la capacità di codifica può scendere sensibilmente. Supponendo uno scenario cautelativo con una capacità di 5 bit per token (con token mediamente composti da 4 caratteri), per codificare i 696 bit del messaggio Alice dovrà generare un testo di circa 140 token (circa 560 caratteri). Tale lunghezza è perfettamente compatibile con una normale risposta di un assistente virtuale.

Per segnalare la fine della comunicazione, può essere utilizzato un carattere speciale che indichi al destinatario il termine del messaggio, per quanto non sia fondamentale in quanto, in caso di assenza di questo carattere, la decodifica procederà normalmente con l'unica differenza che i token decodificati nella fase finale del messaggio sembreranno casuali.

L'output che Eve osserverà sarà simile al seguente:
\begin{lstlisting}[caption={Esempio di testo generato da Meteor contenente un messaggio cifrato. Il messaggio è in formato Markdown, il principale stile di formattazione degli LLM moderni.}, captionpos=b, label={lst:meteor-output-example}]
"Il dominio di una funzione di due variabili e' l'insieme delle coppie $(x, y) \in \mathbb{R}^2$ per le quali la funzione e' definita. Consideriamo ad esempio $f(x, y) = \sqrt{1 - x^2 - y^2}$; affinche' la radice quadrata sia definita, deve valere $1 - x^2 - y^2 \ge 0$, ovvero $x^2 + y^2 \le 1$, che rappresenta il cerchio unitario chiuso nel piano cartesiano. Prendiamo invece $g(x, y) = \frac{1}{x^2 + y^2 - 1}$: il denominatore si annulla quando $x^2 + y^2 = 1$, quindi il dominio e' tutto $\mathbb{R}^2$ esclusa la circonferenza di raggio~1. Per $h(x, y) = \log(x - y)$, la condizione e' $x - y > 0$, ossia il semipiano al di sopra della bisettrice $y = x$. In generale, per determinare il dominio di una funzione multivariata, dobbiamo identificare tutte le condizioni che potrebbero rendere l'espressione non definita, come radici di indice pari di argomenti negativi, denominatori nulli o argomenti di logaritmi non positivi. La soluzione e' tipicamente un sottoinsieme di $\mathbb{R}^2$ o $\mathbb{R}^3$, la cui rappresentazione grafica risulta spesso essenziale per visualizzare le regioni di validita' della funzione."
\end{lstlisting}

In questo testo, i token non sono stati scelti tramite un campionamento casuale (\emph{random sampling}) sulla testa della distribuzione, ma tramite la codifica dei bit del messaggio segreto. Per ogni passo, \emph{Meteor} ha generato la distribuzione di probabilità tramite il modello linguistico, ha rimosso la coda della distribuzione sotto una soglia prefissata e ha ordinato i \emph{logits} ottenuti. Gli indici, distribuiti proporzionalmente alle probabilità dei token in un intervallo da $0$ a $2^p$ (dove $p$ è la precisione), sono stati usati per selezionare il token i cui bit più significativi corrispondessero ai bit successivi del messaggio cifrato.

Se, ad esempio, i bit da codificare fossero \texttt{0101110101...} e le distribuzioni di probabilità dei token disponibili coprissero gli intervalli \texttt{0001...}, \texttt{0100...}, \texttt{0110...} e \texttt{1010...}, \emph{Meteor} selezionerebbe il token associato all'intervallo $[\texttt{0100...}, \texttt{0110...})$, codificando i primi due bit (\texttt{01}). Al passaggio successivo, il token selezionato entra a far parte del contesto del modello e i bit codificati vengono rimossi dalla testa del messaggio segreto, procedendo iterativamente fino al completamento o al raggiungimento dei limiti impostati.

Il testo risultante, coerente e naturale, transita sulla rete domestica senza destare sospetti. Bob, ricevendo il messaggio, utilizza lo stesso modello e la medesima chiave per estrarre i bit e decifrare la comunicazione, comprendendo la richiesta di Alice. Grazie a \emph{Meteor}, Alice è riuscita a eludere la sorveglianza di Eve, rendendo la comunicazione indistinguibile da un comune output di un modello linguistico (e magari è riuscita anche ad ottenere la risposta al problema di analisi che stava studiando, ma questo non dipende da \textit{Meteor}).


\section[VQ-GAN]{VQ-GAN: generazione di immagini ad alta risoluzione}

Successivamente, analizzeremo le ragioni per cui architetture alternative, come le VQ-GAN (Vector Quantized Generative Adversarial Networks), sono state considerate per l'estensione di tali concetti al dominio visuale.

\section{Definizioni e terminologia}

Nel seguito sono presentati un insieme di termini e definizioni fondamentali per la comprensione dei concetti chiave trattati nella tesi. I termini sono organizzati per aree tematiche omogenee, procedendo dai fondamenti della comunicazione sicura fino agli aspetti più specialistici della generazione steganografica. Questo lessico tecnico costituisce il vocabolario essenziale per navigare efficacemente all'interno del discorso sulla steganografia generativa, sui modelli linguistici, sulle architetture di reti neurali e sui generatori di immagini.

\subsection*{Concetti fondazionali della comunicazione}
\begin{itemize}[label=\textbullet, leftmargin=*, itemsep=0em]
	\item \textbf{Plaintext}: locuzione che designa un'informazione nella sua forma originaria e integralmente leggibile, antecedente a qualsiasi trasformazione crittografica o steganografica. Nel contesto della presente trattazione, il plaintext costituisce il dato in chiaro che si intende proteggere.
	\item \textbf{Ciphertext}: denominazione attribuita al risultato dell'applicazione di un algoritmo di cifratura a un plaintext. Tale rappresentazione, strutturalmente incomprensibile in assenza della corrispondente chiave di decifrazione, costituisce il prodotto finale del processo crittografico.
	\item \textbf{Coverobject}: oggetto multimediale originale, nella sua forma inalterata, che funge da contenitore apparentemente innocuo all'interno del quale si intende occultare un messaggio segreto.
	\item \textbf{Stegoobject}: oggetto risultante dall'operazione di incorporazione del messaggio segreto all'interno del coverobject. La proprietà fondamentale che tale artefatto deve soddisfare è l'indistinguibilità percettiva e statistica dal coverobject originale.
	\item \textbf{Payload}: quantità di informazione segreta, espressa in bit, che può essere efficacemente occultata all'interno di un coverobject senza comprometterne l'integrità statistica e percettiva.
\end{itemize}

\subsection*{Discipline della comunicazione sicura}
\begin{itemize}[label=\textbullet, leftmargin=*, itemsep=0em]
	\item \textbf{Crittografia}: disciplina formale che si occupa di rendere non intelligibile un messaggio ad attori non autorizzati mediante l'applicazione di metodi di cifratura. La crittografia costituisce la tecnica preminente per garantire la confidenzialità dei dati nei sistemi informatici contemporanei. Tale disciplina si fonda su principi matematici e algoritmi complessi atti a trasformare l'informazione in una forma cifrata, decifrabile esclusivamente da chi possiede la chiave corretta. La crittografia trova applicazione in molteplici ambiti, tra cui la comunicazione sicura, la protezione dei dati personali e la sicurezza delle transazioni telematiche. Pur costituendo il fondamento dei sistemi comunicativi moderni, la crittografia è potenzialmente vulnerabile ad attacchi di forza bruta, a exploit basati su vulnerabilità specifiche degli algoritmi, a compromissioni delle chiavi e, in prospettiva, ad algoritmi di decriptazione quantistica. Cionondimeno, allo stato attuale, quando implementata correttamente, la crittografia fornisce un livello di sicurezza elevato, rendendo computazionalmente proibitivo l'accesso non autorizzato ai dati \cite{neskey-2025-are}.
	\item \textbf{Cifratura simmetrica}: paradigma crittografico in cui il medesimo segreto condiviso viene impiegato tanto per la cifratura quanto per la decifrazione del messaggio. Tale metodologia, generalmente caratterizzata da una maggiore efficienza computazionale rispetto alla cifratura asimmetrica, richiede tuttavia che la chiave sia trasmessa in modo sicuro tra le parti coinvolte nella comunicazione.
	\item \textbf{Steganografia}: disciplina che si occupa di occultare la mera esistenza di un messaggio all'interno di un oggetto multimediale, rendendolo impercettibile a osservatori non autorizzati. La steganografia si configura come tecnica complementare alla crittografia: mentre quest'ultima protegge il contenuto informativo, la steganografia tutela la riservatezza dell'atto comunicativo stesso. Storicamente impiegata per finalità di spionaggio e comunicazione clandestina, la steganografia ha trovato nuove applicazioni nell'era digitale, quali la protezione dei diritti d'autore, la comunicazione sicura e la salvaguardia della privacy in rete. Le modalità di implementazione includono l'inserimento di dati occulti in immagini, tracce audio o flussi video, nonché l'impiego di tecniche di codifica testuale. La steganografia è tuttavia vulnerabile a tecniche di rilevamento avanzate (steganalisi) che analizzano le proprietà statistiche degli oggetti multimediali per identificare anomalie indicative della presenza di dati nascosti.
	\item \textbf{Steganalisi}: disciplina deputata al rilevamento di messaggi occulti all'interno di oggetti multimediali, avvalendosi di metodologie di analisi statistica e di apprendimento automatico. La steganalisi costituisce il corrispettivo avversario della steganografia, concentrandosi sull'individuazione dell'esistenza stessa del messaggio celato piuttosto che sulla sua protezione. Con l'avvento delle tecnologie digitali, la steganalisi ha raggiunto livelli di sofisticazione notevoli, impiegando algoritmi avanzati per l'esame delle caratteristiche statistiche degli oggetti multimediali. Cionondimeno, la sua efficacia è limitata dalla capacità dei modelli di rilevamento di discriminare tra dati naturali e dati generati da modelli linguistici avanzati, quali quelli impiegati in \emph{Meteor}.
\end{itemize}

\subsection*{Fondamenti di modellazione linguistica}
\begin{itemize}[label=\textbullet, leftmargin=*, itemsep=0em]
	\item \textbf{Modelli linguistici}: classe di modelli di intelligenza artificiale progettati per comprendere e generare testo in linguaggio naturale. Tali modelli vengono addestrati su ingenti quantità di dati testuali e impiegano tecniche di apprendimento automatico per apprendere le regolarità grammaticali, il significato delle parole e le relazioni sintattico-semantiche intercorrenti tra le frasi. I modelli linguistici trovano ampia applicazione in sistemi quali chatbot, piattaforme di traduzione automatica, generatori di testo e analizzatori del sentiment. Essi possono tuttavia presentare vulnerabilità riconducibili a bias presenti nei dati di addestramento e a limitazioni nella comprensione del contesto, con possibili ripercussioni sull'accuratezza e sull'appropriatezza dei risultati prodotti.
	\item \textbf{Modelli di linguaggio di grandi dimensioni (LLM)}: sottocategoria di modelli linguistici caratterizzata da un numero di parametri considerevolmente elevato, che consente di apprendere rappresentazioni più complesse del linguaggio naturale e di generare testo caratterizzato da maggiore coerenza e naturalezza. Gli LLM vengono addestrati su corpora testuali di dimensioni massive e si avvalgono di architetture avanzate, quali i Transformer, per catturare le dipendenze a lungo termine tra le unità linguistiche. Le loro applicazioni spaziano dalla generazione testuale alla traduzione automatica, dalla risposta a domande all'analisi del sentiment.
	\item \textbf{Modelli generativi}: classe di modelli di intelligenza artificiale atti a generare nuovi dati che presentano caratteristiche analoghe a quelli impiegati in fase di addestramento. Attraverso tecniche di apprendimento automatico, tali modelli apprendono la distribuzione sottostante ai dati di training e producono nuovi campioni che seguono la medesima distribuzione. I modelli generativi trovano impiego in numerosi ambiti, inclusi la generazione testuale, la sintesi di immagini, la composizione musicale e la sintesi vocale.
\end{itemize}

\subsection*{Componenti dei modelli linguistici}
\begin{itemize}[label=\textbullet, leftmargin=*, itemsep=0em]
	\item \textbf{Token}: unità testuale elementare impiegata nei modelli linguistici, corrispondente a una parola, un carattere o una sottoporzione di parola. I token costituiscono l'unità fondamentale su cui operano i modelli linguistici, rappresentando le entità semantiche minime che il modello deve elaborare e generare. La tokenizzazione — il processo di suddivisione del testo in token — può essere effettuata secondo diverse modalità, quali la segmentazione per parole, per caratteri o per sottoparole (ad esempio mediante algoritmi \emph{subword} come Byte Pair Encoding). La strategia di tokenizzazione adottata influenza la capacità del modello di comprendere e generare testo in modo efficace.
	\item \textbf{Vocabolario}: insieme dei token che un modello linguistico è in grado di processare e generare. Il vocabolario viene definito durante la fase di addestramento e la sua dimensione può variare considerevolmente in funzione del modello e dell'applicazione. Un vocabolario più esteso consente al modello di coprire una gamma più ampia di espressioni linguistiche, ma può comportare un aumento della complessità computazionale e dei tempi di addestramento.
	\item \textbf{Logits}: valori numerici prodotti dall'ultimo strato di un modello linguistico prima dell'applicazione della funzione di attivazione (tipicamente la softmax) per ottenere una distribuzione di probabilità sui token successivi. I logits rappresentano le preferenze del modello per ciascun token del vocabolario, dove valori più elevati indicano una maggiore propensione alla selezione. La funzione di attivazione converte tali valori in probabilità utilizzabili per il campionamento o la selezione deterministica del token successivo.
	\item \textbf{Softmax}: funzione di attivazione impiegata nei modelli linguistici per trasformare il vettore dei logits in una distribuzione di probabilità sull'insieme dei token del vocabolario. Formalmente, la funzione softmax è definita come $\mathit{softmax}(x_i) = e^{x_i} / \sum_{j} e^{x_j}$, dove $x_i$ rappresenta il logit corrispondente al token $i$-esimo e la somma al denominatore è estesa all'intero vocabolario. Il vettore risultante esprime, per ciascun token, la probabilità che esso venga selezionato come elemento successivo della sequenza.
\end{itemize}

\subsection*{Meccanismi di generazione e controllo}
\begin{itemize}[label=\textbullet, leftmargin=*, itemsep=0em]
	\item \textbf{Prompt}: sequenza testuale iniziale fornita in input a un modello linguistico al fine di orientare e condizionare la generazione del testo successivo, definendo il contesto semantico e tematico dell'output prodotto.
	\item \textbf{Contesto}: insieme di token già generati che vengono impiegati come input per il modello linguistico al fine di produrre il token successivo, costituendo la memoria sequenziale su cui si fonda il processo generativo autoregressivo.
	\item \textbf{Entropia}: misura della quantità di informazione contenuta in un messaggio o in una distribuzione di probabilità, formalizzata nell'ambito della teoria dell'informazione. Un messaggio caratterizzato da elevata entropia risulta maggiormente imprevedibile e, pertanto, più adatto ad applicazioni steganografiche, in quanto offre una maggiore ridondanza statistica in cui occultare informazioni senza introdurre anomalie rilevabili.
	\item \textbf{Campionamento}: processo di selezione di un token successivo a partire da una distribuzione di probabilità generata da un modello linguistico. Le modalità di campionamento includono il campionamento casuale (\emph{random sampling}), la selezione deterministica del token più probabile (\emph{greedy decoding}) e il campionamento vincolato, come nel caso di \emph{Meteor}, in cui la scelta del token è determinata dai bit del messaggio segreto.
\end{itemize}

\subsection*{Architetture di reti neurali}
\begin{itemize}[label=\textbullet, leftmargin=*, itemsep=0em]
	\item \textbf{Layer}: strato funzionale di una rete neurale, costituente l'unità architetturale fondamentale delle architetture di deep learning. Ciascun layer svolge trasformazioni specifiche sui dati in ingresso, quali l'estrazione di caratteristiche, l'apprendimento di rappresentazioni intermedie o la generazione dell'output. Le principali tipologie includono:
	\item \textbf{Transformer}: architettura di rete neurale basata sul meccanismo di attenzione (\emph{self-attention}), ampiamente adottata nei modelli linguistici di grandi dimensioni per la sua capacità di catturare relazioni a lungo termine tra gli elementi di una sequenza, superando le limitazioni delle architetture ricorrenti nella modellazione di dipendenze distanti.
	\item \textbf{Convolutional Layer}: strato che applica operazioni di convoluzione per estrarre caratteristiche locali dai dati in ingresso, comunemente impiegato nelle reti neurali convoluzionali (CNN) per l'elaborazione di segnali strutturati su griglia, quali immagini e volumi multidimensionali.
	\item \textbf{Recurrent Layer}: strato che introduce connessioni ricorrenti per mantenere e propagare informazioni sullo stato precedente della sequenza, tipicamente utilizzato nelle reti neurali ricorrenti (RNN) per l'elaborazione di dati sequenziali, quali testo, serie temporali e segnali audio.
	\item \textbf{Adversarial Layer}: strato impiegato nei modelli generativi avversari (GAN) con funzione discriminativa, deputato a distinguere tra campioni reali e campioni sintetici prodotti dal generatore, contribuendo al miglioramento progressivo della qualità dei dati generati attraverso un processo competitivo.
\end{itemize}

\subsection*{Architetture per la generazione visuale}
\begin{itemize}[label=\textbullet, leftmargin=*, itemsep=0em]
	\item \textbf{Generative Adversarial Network (GAN)}: architettura di rete neurale composta da due componenti distinti — un generatore e un discriminatore — che competono tra loro durante la fase di addestramento. Il generatore tenta di produrre dati sintetici indistinguibili da quelli reali, mentre il discriminatore si adopera per discriminare tra campioni autentici e campioni generati. Questo processo di apprendimento avversario conduce a un miglioramento progressivo della fedeltà e della qualità dei dati sintetici prodotti.
	\item \textbf{Vector Quantized-Generative Adversarial Network (VQ-GAN)}: architettura di rete neurale che integra i principi dei GAN con tecniche di quantizzazione vettoriale, impiegando un codebook discreto per rappresentare le caratteristiche latenti dei dati in ingresso. La VQ-GAN è progettata per migliorare la qualità e la stabilità del processo generativo, riducendo al contempo la complessità computazionale mediante la discretizzazione dello spazio latente.
	\item \textbf{Codebook}: insieme strutturato di vettori latenti appresi durante l'addestramento di un modello basato su quantizzazione vettoriale, come la VQ-GAN. Ciascun vettore del codebook costituisce un elemento atomico del dizionario di rappresentazioni discrete mediante il quale i dati originali vengono codificati e successivamente ricostruiti.
	\item \textbf{Quantizzazione}: processo di discretizzazione di uno spazio di rappresentazione continuo in un insieme finito di valori discreti. Tale procedura viene impiegata in architetture come la VQ-GAN al fine di ridurre la complessità del modello e migliorare la stabilità del processo di addestramento, trasformando lo spazio latente continuo in un repertorio finito di codici.
\end{itemize}

\subsection*{Misure e garanzie nei sistemi steganografici}
\begin{itemize}[label=\textbullet, leftmargin=*, itemsep=0em]
	\item \textbf{Bitrate variabile}: strategia di codifica in cui la quantità di informazione incorporata in ciascun token può variare dinamicamente in funzione delle caratteristiche del contesto e della distribuzione di probabilità dei token successivi. In contesti a bassa entropia, il sistema può astenersi dalla codifica di informazione, mentre in scenari ad alta entropia può codificare un numero maggiore di bit per token, massimizzando così il throughput informativo senza compromettere la naturalezza statistica del testo generato.
	\item \textbf{Codifica adattiva}: approccio alla codifica steganografica in cui la strategia di inserimento dell'informazione segreta si adatta dinamicamente alle proprietà del contesto e alla distribuzione di probabilità dei token successivi. Tale metodologia consente di massimizzare il throughput informativo nei contesti ad alta entropia, preservando al contempo l'integrità statistica del testo in quelli a bassa entropia, garantendo in tal modo un livello di sicurezza steganografica ottimale.
	\item \textbf{Throughput informativo}: quantità di informazione segreta, misurata in bit, che può essere codificata e trasmessa attraverso un canale di comunicazione, quale un testo generato da un modello linguistico. Il throughput informativo è determinato da molteplici fattori, tra cui la capacità di codifica del modello sottostante, la distribuzione di probabilità dei token successivi e la strategia di codifica adottata (bitrate variabile o codifica adattiva). La massimizzazione del throughput informativo costituisce un obiettivo fondamentale nella progettazione di sistemi di steganografia generativa, in quanto consente di veicolare una maggiore quantità di informazione occulta senza pregiudicare la naturalezza del testo prodotto.
	\item \textbf{Integrità statistica}: proprietà di un testo generato da un modello linguistico in virtù della quale le sue caratteristiche statistiche — quali la distribuzione dei token, la frequenza delle parole e la coerenza sintattica — risultano coerenti con quelle di un testo naturale non steganografico. Il mantenimento dell'integrità statistica è essenziale per prevenire il rilevamento da parte di strumenti di steganalisi, costituendo un requisito fondamentale per l'efficacia dell'occultamento.
	\item \textbf{Anomalia semantica}: incoerenza o incongruenza nel significato di un testo generato, riconducibile a selezioni di token non naturali o a errori nel processo generativo. Le anomalie semantiche costituiscono indicatori potenzialmente rivelatori per strumenti di steganalisi, in quanto segnalano una deviazione dal comportamento generativo naturale. L'elusione di tali anomalie è pertanto cruciale per garantire la sicurezza steganografica del sistema.
	\item \textbf{Soglia di entropia minima}: valore critico del parametro entropia al di sotto del quale un sistema di steganografia generativa, quale \emph{Meteor}, si astiene dalla codifica di informazione segreta. Tale soglia viene impiegata per garantire che l'inserimento di informazioni occulte avvenga esclusivamente in contesti in cui la distribuzione di probabilità dei token successivi presenti una varianza sufficiente, prevenendo in tal modo la generazione di anomalie semantiche potenzialmente rilevabili.
	\item \textbf{Soglia di probabilità}: valore prefissato impiegato per troncare la distribuzione di probabilità dei token successivi prodotta da un modello linguistico, eliminando i token la cui probabilità risulta inferiore a tale soglia. Questo meccanismo, noto come troncamento della coda (\emph{tail truncation}), consente di circoscrivere le scelte ai token maggiormente probabili, migliorando la coerenza e la naturalezza del testo generato e mitigando il rischio di introdurre anomalie semantiche o statistiche.
	\item \textbf{Troncamento della coda}: tecnica impiegata nei modelli linguistici per limitare l'insieme dei token candidati alla selezione a quelli la cui probabilità supera una soglia prefissata. Tale procedura migliora la coerenza e la naturalezza del testo generato, eliminando i token meno probabili che potrebbero introdurre anomalie semantiche o statistiche. Questa tecnica riveste particolare rilevanza in sistemi di steganografia generativa, come \emph{Meteor}, in quanto contribuisce a garantire che il testo prodotto rimanga statisticamente indistinguibile da un testo naturale, anche in presenza di informazione occulta codificata al suo interno.
\end{itemize}